% 結果セクションでは、測定データを表やグラフで示す
% 表とグラフの例を以下に示す

\section{結果}

% サブセクションを使って結果を分類することを推奨
\subsection{測定結果の概要}
% ここに測定結果の概要を記述します。

% ===== 表の例 =====
% 表はtable環境とtabular環境を使用
% [H]オプションでその場所に配置（floatパッケージが必要）
% \SI{}コマンドで単位を記載（siunitxパッケージが必要）
% 【例】負荷抵抗特性の表:
% \begin{table}[H]
%   \centering
%   \caption{負荷抵抗に対する出力電圧の変化}
%   \label{table:負荷抵抗特性}
%     \begin{tabular}{>{\hspace{4mm}}c<{} >{\hspace{4mm}}c<{} >{\hspace{4mm}}c<{} >{\hspace{4mm}}c<{}}
%       \noalign{\hrule height 1.2pt}
%       負荷$R_L$ & $V_{\mathrm{in}}$ [Vpp] & $V_{\mathrm{out}}$ [Vpp] & 電圧増幅率 \\
%       \hline
%       開放 & 2.02 & 18.3 & 9.06 \\
%       \SI{100}{\kilo\ohm} & 2.04 & 16.6 & 8.14 \\
%       \SI{10}{\kilo\ohm} & 2.04 & 9.36 & 4.59 \\
%       \SI{1}{\kilo\ohm} & 2.03 & 1.69 & 0.830 \\
%       \hline
%     \end{tabular}
% \end{table}

% ===== グラフの例 =====
% 画像はfigure環境で挿入
% \includegraphicsでファイルを指定（graphicxパッケージが必要）
% width=0.7\linewidthで幅を調整
% 【例】波形のグラフ:
% \begin{figure}[H]
%   \centering
%   \includegraphics[width=0.7\linewidth]{wave_graph.png}
%   \caption{負荷抵抗を接続しなかった時の波形のグラフ}
%   \label{fig:波形}
% \end{figure}

% ===== 複数の図を横に並べる例 =====
% minipage環境を使用して複数の図を並べる
% 【例】2つの図を並べる:
% \begin{figure}[H]
%   \centering
%   \begin{minipage}{0.45\textwidth}
%     \centering
%     \includegraphics[width=\linewidth]{image1.jpg}
%     \caption{図1のキャプション}
%     \label{fig:image1}
%   \end{minipage}
%   \hfill
%   \begin{minipage}{0.45\textwidth}
%     \centering
%     \includegraphics[width=\linewidth]{image2.jpg}
%     \caption{図2のキャプション}
%     \label{fig:image2}
%   \end{minipage}
% \end{figure}

% ===== グラフへの言及の例 =====
% 図\ref{fig:ラベル名}や表\ref{table:ラベル名}で参照
% 【例】
% 表\ref{table:負荷抵抗特性}に負荷抵抗特性の測定結果を、
% 図\ref{fig:波形}に波形のグラフを示す。
% 図\ref{fig:電圧増幅率}は周波数と電圧増幅率を対数でとった両対数グラフである。
% 最小二乗法を用いて近似曲線を求めた結果、決定係数は $R^2 = 0.990$ となった。
